\par\medskip\noindent
\underline{Corolario} 1.3.6. \underline{Supongamos que todo
morfismo de haces de conjuntos} $P\times G \longrightarrow G$
\underline{que sea trivial sobre} $P\times e_G$ \underline{y sobre}
$\sisi{e_P}{e_P}\times G$ \underline{es trivial}. \underline{Entonces},
\underline{si} $G$ \underline{es conmutativo}, $G$ \underline{es
central en toda extensión de grupos de} $P$ \underline{por} $G$.
\underline{Si} $P$ \underline{es también conmutativo}, \underline{y
si se satisface la hipótesis de} 1.3.5 a), \underline{entonces toda
extensión de grupos} $E$ \underline{de} $P$ \underline{por} $G$
\underline{es conmutativa}.

Para la primera afirmación hay que demostrar que la operación de $P$
sobre $G$ definida por la extensión $\sisi{G}{E}$ es trivial. Ahora
bien, esta operación se expresa mediante un morfismo
$u : P\times G \longrightarrow G$, tal que las restricciones a
$P\times e_G$ y a $\sisi{e_P}{e_P}\times G$ del morfismo
correspondiente $(p,g) \longmapsto u(p,g)g^{-1}$ son evidentemente
triviales; por tanto, este último morfismo es trivial, es decir, se
tiene $u(p,g) = g$, que es lo que se quería demostrar. Observemos ahora
que, si $E$ es una extensión de $P$ por $G$, entonces el grupo opuesto
$E^o$ a $E$, que es una extensión de $P^o$ por $G^o$, tiene como
estructura de bitorsor subyacente $i(E^o)$ la deducida del bitorsor
$i(E)$ intercambiando en él las operaciones izquierda y derecha. Cuando
$E$ es una extensión central, se tiene, pues, $i(E)=i(E^o)$, y si,
además, $P$ es conmutativo, es decir, $P=P^o$, la afirmación de unicidad
1.3.5 a) da $E=E^o$, es decir, $E$ es conmutativo. Esto demuestra el
corolario.


%%%%%%%%%%  scan idx 163  |  folio 152  %%%%%%%%%%

1.3.7. Un caso particular interesante en el que se satisfacen todas
las hipótesis consideradas en 1.3.5 y 1.3.6 es aquel en que se tiene

\begin{equation}\tag{1.3.7.1}
\underline{\mathrm{Hom}}_{pt}(P,G) \;=\; e \ \text{(haz final)} \qquad ,
\end{equation}
donde el primer miembro designa el haz de los morfismos de haces de
conjuntos apuntados de $P$ en $G$. Esta condición se escribe también,
independientemente de las secciones marcadas, en la forma (análoga a
(1.3.1.2))

\begin{equation}\tag{1.3.7.2}
G \;\simeq\; \underline{\mathrm{Hom}}_{ens}(P,G) \qquad ,
\end{equation}
donde el segundo miembro designa el haz de los morfismos de los haces
de conjuntos subyacentes a $P$ y $G$. Estas condiciones significan
también que, para todo objeto $Q$ de $\underline{T}$, la aplicación
$u \longmapsto u\circ pr_2$

\begin{equation}\tag{1.3.7.3}
\mathrm{Hom}_{ens}(Q,G) \;\xrightarrow{\ \sim\ }\; \mathrm{Hom}_{ens}(P\times Q,G)
\end{equation}
es biyectiva. Se ve que esta propiedad de $P$ (para $G$ fijo) es estable
por productos cartesianos, y por tanto implica, en particular, la misma
propiedad para los productos cartesianos $P\times P$,
$P\times P\times P$, etc. Es entonces inmediato que implica las
hipótesis formuladas en 1.3.5 y en 1.3.6; además, implica que las
categorías que intervienen en (1.3.4.1) son todas \underline{rígidas},
es decir, que todo automorfismo de un objeto de una de estas categorías
es la identidad. Además, como la condición (1.3.7.1) sobre $P,G$ es
estable por todo cambio de base $S \longrightarrow e$, se deduce de
inmediato por descenso, a partir del resultado de unicidad 1.3.5 a),
que, para que un bitorsor rigidificado $E$ sobre $P$ corresponda en
efecto a una extensión de $P$ por $G$, basta que así sea
\underline{localmente sobre} $e$, lo cual (en virtud de 1.3.5 b)) se
expresa también mediante la anulación de la sección

\begin{equation*}
\underline{\delta}(E) \in \Gamma\, R^1 h_*(G_{P\times P}) 
\end{equation*}

%%%%%%%%%%  scan idx 164  |  folio 153  %%%%%%%%%%


definida por $\delta(E)$, donde

% Canon: SGA-SRC-001259.
\begin{equation*}
h : \sisi{P\times P}{P\times P} \longrightarrow e
\end{equation*}
es el morfismo estructural.
